% !TeX root = ../papers/nc-p3.tex
\section{七十一節点の双三次式による十五次元成分の局所分類}
\label{sec:marked-bicubic-nodes}
\begin{theorem}\label{thm:marked-bicubic-nodes}
十五次元成分上の印付き被覆のうち、例 $t$ の素数2657での第一位置の印
$(1,1,0,0)$ に達する既約成分を取る。その一般の窓で
\[
 D=\{\det[Q_l,P_r]=0\}\subset\mathbb P^2_l\times\mathbb P^2_r,\qquad
 D'=\{\det[G_x,F_z]=0\}\subset\mathbb P^2_x\times\mathbb P^2_z
\]
はともに71個の通常二重点だけを特異点とする双次数 $(3,3)$ の超曲面である。
各々36点では行列階数が1である。残る35点は二方向で標準的に対応し、
十九点の窓の点列と、二つの互いに素な八点集合に分かれる。
特に定理\ref{thm:marked-cubic-surface-pairs}の残余八点はこの一般点で被約である。

$D$ の71節点を保つ双三次式の被約局所成分は滑らかで、式の全体倍率を含む次元は32、
両平面の独立な射影変更と倍率を除いた一般次元は15である。
印付き窓からこの成分の商への写像は支配的かつ生成的有限で、一般点で分離的な局所座標を与える。
従って一般の局所変形は、この節点数を保つ双三次式の変形で追跡できる。
\end{theorem}

\begin{proof}[二方向の行列と節点]
四線型形式を $f(l,r,x,z)=x[Q_l,P_r]z=-l[G_x,F_z]r$ と置く。
四つの文字について全ての一階偏微分が零となる四つ組を考える。
$[Q_l,P_r]$ の階数が2ならその左右核 $x,z$ は一意で、
$\det[Q_l,P_r]$ の微分は余因子公式によって $f$ の $l,r$ 微分の非零倍率である。
従って $D$ の階数2の特異点はこのような四つ組に持ち上がる。
通常節点しかないときには、四つ組のいずれの行列も階数1以下にならない。
実際、階数0なら行列式が三重以上になる。
階数1なら単元を消した二行二列のSchur補の行列式が局所式である。
$x,z$ はその左右核の二空間に属し、$x\,d[Q_l,P_r]\,z=0$ が全方向で成立する。
これは四成分の一次部の非自明な線型関係であり、二次主要部は退化して通常節点でない。
もう一方についても同じなので、二つの階数2の特異点集合が左右核によって一対一に対応する。

定理\ref{thm:marked-cubic-surface-pairs}の二つの曲面対の交叉からこの四重臨界集合に点が入る。
例えば第一対で $G_x,F_z$ はいずれも $l,r$ を左右の同一固有値の固有ベクトルとして持つ。
従って $l[G_x,F_z]=[G_x,F_z]r=0$ であり、他方の二式は
$l[G_x,F_{z'}]r=l[G_{x'},F_z]r=0$ から従う。これは交点スキーム上の等式でもある。
どちらの交叉も対応する双三次式の特異スキームに含まれ、通常節点だけなら長さ27の交叉は被約である。
両方向の27点が同じ四つ組になる条件は、$Q_l,P_r$ についても $x,z$ がそれぞれ左右の
同一固有値の組になることである。四つの固有値を星座標に戻せば四辺の点対の全式が成り立つ。
逆も明らかで、重なるのは星点以外の19本の点列に限る。従って階数2の節点は少なくとも35点である。

一般点では交叉交換子は跡零八空間の同型である。
$[Q_l,P_r]$ の射影行列はSegre四次元多様体を可逆行列の点から跡零空間へ射影し、
線型同型を施したものになる。射影の中心は割線多様体の外なのでその像は次数6である。
階数1・跡零の行列の多様体はもう一つのSegreの超平面切断で次元3、次数6である。
$\mathbb P^7$ での交叉が有限なら長さ36、通常節点ではSchur補の四一次部が独立なので単純である。
よって節点の上界71と通常性が成立すれば、二方向に全ての点が上の記述で尽くされる。
\end{proof}

\begin{proof}[上界、節点条件の接空間と支配性]
上の印の剰余点で両行列式を直接計算する。各平面の先頭非零座標を1とした九つの座標片で、
元の四アフィン方向の微分を全て課した後に先行座標を零にする。
二方向とも特異点の長さは71、二次小行列式を加えた長さは36、
四行Hesse行列式も零とした軌跡は空である。
特異点の射影性と非退化Hesse条件から、局所エタールに各点では四偏微分を解いた
唯一の臨界点列があり、その点での式の値の消滅が節点が残る条件である。
従ってどちらも近傍で特異点は高々71、全て通常節点である。
標数零の一般点に前段の下界を適用すれば等号であり、71個全てが残る。
このため $t$ の該当する標数零の持上げにもその全てが特殊化する。

100次元の双三次式のアフィン空間で、節点存続式の微分は各節点での変分三次式の評価である。
$D$ の剰余の全節点環に100個の単項式を写した階数は68である。
従って七十一節点を保つ局所軌跡の接次元は32以下である。
一方、四つの跡零三行列組の96係数について交換子方程式の微分核は34次元である。
この印付き正規化表示が滑らかな34次元であることは、$t$ で滑らかな76次元の
二切断軌跡、単純点のエタールな選択、二つの反対行列と跡の規格化から従う
（定理\ref{thm:marked-mixed-bracket-reconstruction}と同じ枠の計数）。
剰余でも微分核が同次元なので接方向の基底を持ち上げられる。
これらの方向で $D$ の変分を計算し、さらに二平面の独立な基底変更と式の倍率を加える。
軌道方向だけで17、窓の変分だけで24、合わせて32個の一次独立な式が得られる。
従って窓と座標変更の像の次元は少なくとも32である。近傍ではその像が全節点の存続式を満たすので
先の接上界と一致し、局所環は正則、像はこの成分に稠密となる。
射影同値を除いた次元は $32-17=15$、元の印付き商と一致する。
\path{computations/verify_marked_bicubic_nodes.py} は両特異イデアル、評価行列、
および行列式の余因子微分を検算する。変分多項式の独立性の下界には130組の代入値を用いる。
\end{proof}
